% ============================================================
%  第二部分之四：前沿方法（2024--2026）
% ============================================================

\section{前沿：2024--2026 的方法演进主线}
\label{sec:frontier}

前两章的方法骨架在 2010 年前后已经成型。2024--2026 年的研究进展并不是推翻了 HMM 或变点检测，而是围绕三个被长期忽视的工程短板展开修补：\emph{持续性的显式控制}、\emph{状态身份在滚动重估中的稳定性}、以及\emph{交易成本与执行延迟的现实约束}。本章沿四条主线梳理最新进展，并在最后给出一个批判性的综合判断。

\subsection{主线一：把“持续性”显式化——统计跳变模型}

HMM 的持久性由转移矩阵隐含表达（$a_{ii}$ 多大、状态就多黏），这带来两个问题：其一，持久性无法\emph{直接}调整，只能间接影响；其二，$\lambda$ 的“漂移—跳变”连续谱在 HMM 中不可表达（其状态切换的概率是常数，与数据无关）。

统计跳变模型（Statistical Jump Model, JM）~\cite{shu2024a,shu2024b} 把持久性提为主体：直接对状态路径求解
\begin{equation}
\min_{s_{1:T} \in \{1,\dots,K\}^T}
\;\; \sum_{t=1}^{T} \mathrm{cost}\big(Y_t,\, \theta_{s_t}\big)
\;+\; \lambda \sum_{t=2}^{T} \mathbb{1}\big[s_t \neq s_{t-1}\big],
\label{eq:jm}
\end{equation}
其中第一项是拟合代价（聚类型：状态参数 $\theta_k$ 由该状态样本的均值等统计量给出），第二项以显式价格 $\lambda$ 惩罚每一次跳变。

式~\eqref{eq:jm} 揭示了一个被 HMM 表述掩盖的连续谱：\textbf{$\lambda = 0$ 时精确退化为 K-means}（逐点聚类、无时序约束）；$\lambda$ 增大则跳变越来越“昂贵”，状态路径越来越平滑；$\lambda \to \infty$ 时所有样本坍缩到单一状态。HMM、K-means 与“跳变控制”由此被统一为同一目标的三个极限。

\begin{table}[htbp]
\centering
\footnotesize
\begin{tabular}{>{\raggedright\arraybackslash}p{2.2cm}>{\raggedright\arraybackslash}p{5.0cm}>{\raggedright\arraybackslash}p{4.6cm}}
\toprule
& \textbf{HMM} & \textbf{统计跳变模型（JM）} \\
\midrule
持久性来源 & 转移矩阵（隐含） & 跳变惩罚 $\lambda$（显式） \\
发射假设 & 概率分布（如高斯） & 仅需状态统计量（均值型） \\
参数估计 & EM（似然） & 动态规划/坐标下降（代价最小化） \\
超参数选择 & 状态数 $K$ & $K$ 与 $\lambda$（可用策略表现做时序 CV） \\
失效模式 & 状态抖动、局部最优 & $\lambda$ 过大导致状态坍缩 \\
\bottomrule
\end{tabular}
\caption{HMM 与统计跳变模型的结构对比}
\label{tab:hmm-vs-jm}
\end{table}

Shu et al.~\cite{shu2024a} 的实证给出了这套方法的标准评估范式：在美、德、日主要股指 1990--2023 样本上，特征只取收益序列衍生的风险与收益度量；跳变惩罚经时间序列交叉验证选择（直接优化策略表现而非似然）；全部结果含交易成本与 1 日交易延迟。结论：JM 引导的策略在波动率与最大回撤上系统性低于 HMM 引导策略与买入持有，风险调整收益更高；且对交易延迟的鲁棒性更强（延迟 1--2 日的绩效衰减明显小于 HMM 版本）。后者正是“显式持续性”的直接工程回报——状态切换频率低，执行延迟的冲击自然更小。

\subsection{主线二：身份稳定性——Wasserstein HMM 与模板追踪}

第二章的 label switching 问题在滚动重估中有一个时变版本：每次重估后，状态编号都可能重排或含义漂移。Boukardagha~\cite{boukardagha2026} 的 Wasserstein HMM 针对这一痛点提出了一套可直接落地的机制：

\begin{enumerate}[itemsep=2pt]
\item \textbf{滚动高斯 HMM 推断}：保持经典结构，逐窗估计；
\item \textbf{预测性模型阶数选择}：允许状态复杂度随数据自适应（而非固定 $K$）；
\item \textbf{模板身份追踪}：以高斯分量之间的 $2$-Wasserstein 距离，把当前窗口的状态与历史“模板”对齐——状态的\emph{经济身份}由其在分布空间中的位置定义，而非编号。
\end{enumerate}

其实证（跨资产日度组合、严格因果设定、含交易成本）报告：相对等权与 SPX 买入持有基准，风险调整收益更高（报告 Sharpe $2.18$ vs $1.59$ 与 $1.18$），最大回撤显著更低（$-5.43\%$ vs $-14.62\%$）；在 2025 年初的股票抛售中动态减配股票、转向防御资产。与同特征的非参数 KNN 对照，参数化 Regime 模型的换手率与权重演化平滑度显著更优。

\begin{remark}[对报告数字的正确读法]
上述数字来自单篇工作论文，应作为“该设计在特定样本上的表现证据”，而非可外推的承诺。真正可迁移的是其\emph{结构性结论}：在日频再平衡的设置下，\emph{状态身份的保持}（identity preservation）与\emph{自适应复杂度控制}是回撤与实现稳健性的\emph{一阶决定因素}——这比参数选择或特征工程更靠前。换句话说，Regime 系统的稳定性问题，本质上是“状态身份是否会漂移”的问题。
\end{remark}

\subsection{主线三：识别与预测分离的两阶段架构}

Shu et al.~\cite{shu2024b} 给出了当前最清晰的模块化范式，把系统切成三段，每段的输入输出与验证方式彼此独立：

\begin{center}
\begin{tabular}{>{\raggedright\arraybackslash}p{3.7cm}>{\raggedright\arraybackslash}p{3.7cm}>{\raggedright\arraybackslash}p{3.7cm}}
\toprule
\textbf{识别}（无监督） & \textbf{预测}（监督） & \textbf{优化}（约束） \\
\midrule
JM 生成历史状态标签（无标签、可解释） & GBDT 以资产特征 $+$ 宏观特征预测下一期状态 & 交易成本感知的均值—方差优化 \\
\bottomrule
\end{tabular}
\end{center}

该架构的方法论价值：识别环节的评估（状态画像是否清晰、是否稳定）与预测环节的评估（分类的样本外准确率、校准）可以\emph{分别}进行，避免了“端到端指标掩盖局部故障”的黑箱困境；优化环节则把交易成本与约束显式化。其 1991--2023 年跨资产组合实证（全球股票、债券、房地产、商品）报告了对最小方差、均值方差与朴素分散基准的一致改善。

\subsection{主线四：约束、成本与实现现实}

识别得再好，若忽略实现摩擦，收益会被成本吃掉。这一主线有两份关键证据：

\textbf{换手与杠杆约束的价值}。Lewin \& Campani~\cite{lewin2023} 在含交易成本的样本外测试中比较了三种 Regime 配置：无约束、低换手控制（LoT，对微小再平衡设置动态阈值）与最大杠杆控制（MaxLev，动态调整风险厌恶参数以约束杠杆）。结论具备直接工程含义：考虑成本后，带 LoT 的模型风险回报指标优于不带者；MaxLev 版本在含成本情形下接近甚至超过无约束基准；两者在确定性等价收益上统计显著地优于单状态模型与无约束 Regime 模型。

\textbf{传染状态下的防御逻辑}。Scherer~\cite{scherer2020} 对另类风险溢价组合的传染分析给出资产配置层面的结论：当策略间传染发生时，分散化失效，最优回应是降低杠杆并回避对传染状态敏感的资产——这为“高波动防御型 Regime $\to$ 系统性降杠杆”提供了组合理论支撑。

\begin{remark}[实现细节不是“工程杂活”]
把两条证据放回第一部分的下界框架：识别精度受信息限制、边际改善有限；而换手、延迟、约束这些\emph{实现变量}的改善空间没有此类统计下界。这解释了为什么 2024--2026 的高质量研究普遍把成本与延迟作为核心实验设计（JM 的延迟鲁棒性测试、Wasserstein HMM 的换手对照、Lewin--Campani 的成本外测），而不是把它们放进“实现细节”一节一笔带过。
\end{remark}

\subsection{补充：混合族与判别式增强}

除四条主线外，还有一类“生成—判别混合”路线值得记录：Koukorinis et al.~\cite{koukorinis2025} 用 HMM 从微观结构时间序列生成特征嵌入，再以核机器（SVM/MKL）做判别，在高频 Regime 分类与交易方向预测上优于单核 SVM、logistic 与前馈网络。其结构含义与第~\ref{sec:engineering}~节体系一致：\emph{生成模型负责把隐结构“特征化”，判别模型负责做最终决策}——这是一种工程上极其可复制的分工：概率模型给证据，判别器给行动。

\subsection{谱系统一视图与批判性综合}

\begin{table}[htbp]
\centering
\footnotesize
\begin{tabular}{p{3.4cm}p{3.0cm}p{3.0cm}p{2.6cm}}
\toprule
\textbf{方法} & \textbf{持续性机制} & \textbf{身份处理} & \textbf{成本/延迟} \\
\midrule
HMM / MS & 转移矩阵（隐含） & 无（编号漂移） & 少见 \\
HSMM / 时长分布 & 显式时长分布 & 无 & 少见 \\
统计跳变模型 & 跳变惩罚 $\lambda$（显式） & 无（可用模板补） & 延迟鲁棒性测试 \\
Wasserstein HMM & 转移矩阵 $+$ 阶数选择 & $2$-Wasserstein 追踪 & 换手对照 \\
两阶段（识别$+$预测） & 由识别模型继承 & 由识别模型继承 & 成本感知优化 \\
约束型配置 & 继承 & 继承 & LoT / MaxLev 显式约束 \\
\bottomrule
\end{tabular}
\caption{前沿方法的能力矩阵：四条主线在每一维度上的贡献}
\label{tab:frontier-matrix}
\end{table}

综合判断：

\begin{enumerate}[itemsep=3pt]
\item 前沿的进步\emph{不是}“预测得更准”。没有任何一项工作声称突破了第一部分的信息论下界；相反，它们都在设计上承认了下界。
\item 前沿的进步在于\textbf{把不可控的统计问题转化为可控的工程问题}：持续性从隐含变显式（可以调、可以验证）、身份从编号变模板（可以追踪、可以审计）、成本从被忽略变前置（可以设计约束、可以度量）。
\item 由此得到一个修正后的判断：\textbf{Regime 系统的价值主要体现在“正确地保持姿态”，而非“更早地预知未来”}。上调乐观程度的地方不在识别精度，而在系统的稳健性、可解释性与成本效率。
\end{enumerate}

\keynote{2024--2026 的前沿没有推翻 HMM，而是给它补了三块短板：持续性先验、身份稳定性、实现现实。方向本身就是一条证据：Regime 系统的第一价值是“正确地保持姿态”，而不是“更早地预知未来”。}